\subsection{Sprint 2}
Para a Sprint 2, o time optou por dividir-se em sub-equipes, cada uma sendo encarregada de uma tarefa diferente. Assim, foram planejadas diversas atividades importantes para o avanço do projeto, como o desenvolvimento da tela de dashboard e todos os seus componentes(dois gráficos, uma tabela e um componente para filtragem dos dados) e sua integração. Essas tasks foram definidas e entregues a uma sub-equipe bem diversa com o objetivo de cumprir os prazos estabelecidos, ao mesmo tempo em que se buscava garantir a qualidade das entregas e permitir que aqueles integrantes que não possuíssem tanto conhecimento técnico pudessem aprender e evoluir durante o desenvolvimento do projeto com o auxílio dos outros colegas de time. Mesmo com a significativa complexidade das atividades previstas e de alguns contratempos enfrentados, todas foram concluídas dentro do período da sprint.

Como mencionado anteriormente, o desenvolvimento não foi isento de desafios. Durante a sprint, um dos membros da sub-equipe, infelizmente, não conseguiu contribuir para o desenvolvimento das tarefas sob sua responsabilidade devido à problemas pessoais. A sua ausência foi comunicada tardiamente, o que dificultou a reorganização das tarefas dentro do curto prazo restante. Para lidar com essa situação, parte das responsabilidades que estavam sob a supervisão desse colega precisaram ser redistribuídas, o que sobrecarregou os demais integrantes da e limitou o tempo que eu tinha disponível para revisar as tarefas dos membros com menos experiência. Essa limitação resultou em alguns componentes que, embora funcionais, apresentaram falhas ao não refletirem integralmente o design do protótipo. Além disso, a falta de tempo também impediu a implementação de testes unitários, algo que acaba comprometendo diretamente a confiabilidade e na robustez do código entregue.

Essa experiência reforçou a importância de uma comunicação aberta e antecipada dentro da equipe. A notificação prévia de problemas, mesmo que pareçam mínimos, é crucial para que o time consiga se reorganizar e mitigar riscos antes que estes se tornem críticos. Além disso, fica evidente que o planejamento deve considerar margens de segurança para lidar com imprevistos, especialmente em sprints com prazos curtos como essa.

Para os próximos passos, o time, como um todo, decidiu que será necessário criar um planejamento que permita padronizar o front-end, assegurando consistência no desenvolvimento de novos componentes e telas. Também será preciso dar continuidade ao trabalho com a implementação da tela de áudio e, caso o tempo permita, realizar os testes unitários que ficaram pendentes nesta sprint. A Sprint 2 trouxe diversos aprendizados valiosos, e, mesmo diante de desafios, o progresso alcançado pela equipe reflete a resiliência e comprometimento do time com os objetivos do projeto.
